\scanpage{55}

% manuscript line(s) 1
Mais le deuxième serait plutôt
% manuscript line(s) 2
\[
\mathrm{Oub}'(L,\Theta)=(L,4\Theta,0)
\]
% manuscript line(s) 2–4
Ce foncteur ne commute \uline{pas} \inserted{néc\textsuperscript{t}} aux \deleted{\unsure{$\otimes$}}\unsure{$\wedge$|$\otimes$} à is près \inserted{(\deleted{\unsure{pas pour|N.B. pas}} $2$ inv. sur $X$,)}), et on en déduit que le foncteur
% manuscript line(s) 2–6
% margin note written sideways
% manuscript line(s) 2–6
\margin{car $(L\otimes L',16\Theta\Theta')$ n'est pas néc. isom. à $(L\otimes L',4\Theta\Theta')$, \deleted{\unsure{puisque pour avoir|p. ex.} $16\Theta\Theta'=0,\ 4\Theta\Theta'=0$} (p.ex. sur \deleted{\unsure{$Z$|$X$}} $\mathrm{Spec}\ \mathbb Z$, $L=L'=\underline O_X$, $\Theta\unsure{=}\Theta'=1$ $\ldots$}
% manuscript line(s) 5
\[
\mathrm{Tors}(\mu_{2/X})\unsure{\mathbin{-}|\to}\underline{\mathrm{Rev.quad}}(X)
\]
% manuscript line(s) 6
ne commute pas aux compositions, contrairement à
% manuscript line(s) 7–8
\diagram{\begin{tikzcd}[ampersand replacement=\&] \underline{\mathrm{Tors}}((\mathbb Z/2\mathbb Z)_X)\arrow[r]\arrow[d]\&\underline{\mathrm{Tors}}(\mu_{2X})\\ \underline{\mathrm{Rev.quad}}(X)\&\end{tikzcd}}

% manuscript line(s) 9
On aura en effet
% manuscript line(s) 10
\[
(L,\delta,T_0)\wedge(L',\delta',T'_0)=(L\otimes L',\delta\cdot\delta',T_0\cdot T'_0)
\]
% manuscript line(s) 11–12
\deleted{\unsure{tant|les}} des \uline{Triples} \uline{étales} $(L,\delta,T_0)$ que des \uline{Triples} (\unsure{quasi|quelconques} étales) $(L,\delta,T_0)$, et
% manuscript line(s) 13
\[
(L,\Theta)\wedge(L',\Theta')=(L\otimes L',\Theta\Theta')
\]
% manuscript line(s) 14
\par\centerline{\rule{1.2cm}{.3pt}}
% manuscript line(s) 14–19
% slanted
% manuscript line(s) 14–19
\margin{Ce qui suit (jusqu'à la page suivante) doit être repris de façon plus fine, en remplaçant \uline{Rev.quad!}$(X)$ par la cat. des $\mu_2$-rev. quadratiques (ou en \unsure{seulement|renonçant|se convaincant} \unsure{algèbres} $\ldots$)}

% manuscript line(s) 15–17
$\underline{\mathrm{Revquad}}(X)$ est une catégorie monoïdale (comm.) par un produit contracté qui prolonge celui sur les $(\mathbb Z/2\mathbb Z)_X$-torseurs

% manuscript line(s) 18–22
$\underline{\mathrm{Revquad}}^{!}(X)$, sous-catégorie pleine formée des rév. quadratiques \inserted{$X'/X$} tels que $T_0=0$ i.e. la trace sur $X'_0/X_0$ $(X_0=V(2.1_X))$ est id. nulle, i.e. $N_{X'_0/X_0}$ est additive i.e. $\varphi_{N_{X'_0/X_0}}$ (forme bil. associée) est id. nulle.

% manuscript line(s) 23
On définit une catégorie monoïdale mixte
% archive sheet number: 54
